\documentclass[12pt, a4paper, oneside]{ctexart}
\usepackage{amsmath, amsthm, amssymb, bm, xcolor, tikz, framed, graphicx, hyperref, mathrsfs}
\usepackage{bookmark}
\usepackage{pgfplots}
\pgfplotsset{compat=1.18}

\newcommand{\bra}[1]{\langle #1 |}
\newcommand{\ket}[1]{| #1 \rangle}
\newcommand{\braket}[2]{\langle #1 | #2 \rangle}

\title{Quantum Mechanics Homework (Week 11)}
\author{Tang Yi \\ [学号已脱敏]}
\usepackage[a4paper, left=2cm, right=2cm, top=2cm, bottom=2cm]{geometry}
\linespread{1.3} % 行距更紧凑
\definecolor{shadecolor}{RGB}{241, 241, 255}

\newcounter{problemname}
\newenvironment{problem}
  {\begin{shaded}\stepcounter{problemname}\par\noindent\textbf{题目 \theproblemname. }}
  {\end{shaded}\par}
\newenvironment{solution}
  {\par\noindent\textbf{题目 \theproblemname 的解答. }}
  {\par}
\newenvironment{note}
  {\par\noindent\textbf{题目 \theproblemname 的注记. }}
  {\par}

\begin{document}

\maketitle

\begin{problem}
  定义径向动量算符：

  \centerline{$p_r = \dfrac{1}{2} \left(\boldsymbol{e}_r \cdot \boldsymbol{p} + \boldsymbol{p} \cdot \boldsymbol{e}_r\right)$}

  其中$\boldsymbol{e}_r = \dfrac{\boldsymbol{r}}{r}$。证明：(1)$p_r$是厄米的；(2)$p_r = -i\hbar\left(\dfrac{\partial}{\partial r} + \dfrac{1}{r}\right)$；
  (3) $[r,p_r] = i\hbar$ ；\\
  (4)$p_r^2 = -\hbar^2 \dfrac{1}{r} \dfrac{\partial^2}{\partial r^2} r$
\end{problem}

\begin{solution}
定义径向动量算符
\[
p_r=\frac12\left(\boldsymbol e_r\cdot\boldsymbol p+\boldsymbol p\cdot\boldsymbol e_r\right),
\qquad
\boldsymbol e_r=\frac{\boldsymbol r}{r},\quad
\boldsymbol p=-i\hbar\nabla.
\]

(1)
由于 $\boldsymbol p^\dagger=\boldsymbol p$ 且 $\boldsymbol e_r^\dagger=\boldsymbol e_r$，有
\[
p_r^\dagger
=\frac12\Big((\boldsymbol e_r\cdot\boldsymbol p)^\dagger+(\boldsymbol p\cdot\boldsymbol e_r)^\dagger\Big)
=\frac12(\boldsymbol p\cdot\boldsymbol e_r+\boldsymbol e_r\cdot\boldsymbol p)
=p_r.
\]
因此 $p_r$ 是厄米的。

(2)
\[
p_r=-i\hbar\left(\frac{\partial}{\partial r}+\frac{1}{r}\right)
\]

作用于 $\psi$：
\[
(\boldsymbol p\cdot\boldsymbol e_r)\psi
=-i\hbar\,\boldsymbol e_r\cdot\nabla\psi
=-i\hbar\,\frac{\partial\psi}{\partial r},
\]
\[
(\boldsymbol e_r\cdot\boldsymbol p)\psi
=-i\hbar\,\nabla\cdot(\boldsymbol e_r\psi)
=-i\hbar\left(\frac{\partial\psi}{\partial r}+(\nabla\cdot\boldsymbol e_r)\psi\right).
\]
三维空间中
\[
\nabla\cdot\boldsymbol e_r=\frac{2}{r}.
\]
代入得
\[
p_r\psi
=-i\hbar\left(\frac{\partial\psi}{\partial r}+\frac1r\psi\right),
\]
即
\[
p_r=-i\hbar\left(\frac{\partial}{\partial r}+\frac1r\right).
\]


(3)
对任意 $\psi$：
\[
rp_r\psi
=-i\hbar(r\partial_r\psi+\psi),
\]
\[
p_r(r\psi)
=-i\hbar\left(\partial_r(r\psi)+\frac1r(r\psi)\right)
=-i\hbar(r\partial_r\psi+2\psi).
\]
故
\[
[r,p_r]\psi=i\hbar\psi,
\quad\Rightarrow\quad
[r,p_r]=i\hbar.
\]

(4)
\[
p_r^2=-\hbar^2\,\frac1r\frac{\partial^2}{\partial r^2}r
\]

因为
\[
p_r=-i\hbar\left(\partial_r+\frac1r\right),
\]
平方得
\[
p_r^2=-\hbar^2(\partial_r+\tfrac1r)(\partial_r+\tfrac1r).
\]
作用于 $\psi$：
\[
(\partial_r+\tfrac1r)(\partial_r+\tfrac1r)\psi
=\partial_r^2\psi+\frac{2}{r}\partial_r\psi.
\]
另一方面，
\[
\frac1r\partial_r^2(r\psi)
=\frac{1}{r}\big(r\partial_r^2\psi+2\partial_r\psi\big)
=\partial_r^2\psi+\frac{2}{r}\partial_r\psi.
\]
因此
\[
p_r^2\psi=-\hbar^2\frac1r\partial_r^2(r\psi),
\]
即
\[
p_r^2=-\hbar^2\frac1r\frac{\partial^2}{\partial r^2}r.
\]
\end{solution}

\begin{problem}
  考虑二维径向波函数在$ r = 0 $附近的渐近解 $R(r)\sim r_a$，束缚态平方可积要求$ a $的取值满足什么条件？
\end{problem}

\begin{solution}
  在二维极坐标下，归一化条件要求
\[
\int_{0}^{\epsilon} |R(r)|^{2}\, r\, \mathrm{d}r < \infty .
\]
若在 $r\to 0$ 处有渐近行为
\[
R(r) \sim r^{a},
\]
则被积函数在 $r=0$ 附近的幂次为
\[
|R(r)|^{2} r \sim r^{2a} r = r^{2a+1}.
\]
积分
\[
\int_{0}^{\epsilon} r^{\,2a+1}\,\mathrm{d}r
\]
收敛的条件是指数满足
\[
2a+1 > -1.
\]
因此得到
\[
\boxed{\,a > -1\,}.
\]
\end{solution}

\begin{problem}
  质量为 $\mu$ 的粒子在如下势场中运动：
\[
V(r)=
\begin{cases}
0, & R_1< r < R_2,\\[4pt]
\infty, & r\le R_1,\ r\ge R_2.
\end{cases}
\]

求其 $l=0$ 的能级和相应的本征波函数。
\end{problem}

\begin{solution}
对于中心势，波函数可分离为球谐与径向部分：
\[
\Psi(\boldsymbol r)=R(r)Y_{lm}(\theta,\phi).
\]
对 \(l=0\)（即 \(Y_{00}=\dfrac{1}{\sqrt{4\pi}}\)）情形，引入径向函数 \(u(r)=rR(r)\)。径向定态薛定谔方程（在势为零的区间 \(R_1<r<R_2\)）化为一维形式：
\[
-\frac{\hbar^2}{2\mu}\frac{d^2u}{dr^2}=E u,\qquad R_1<r<R_2,
\]
并且无穷深势壁给出边界条件
\[
u(R_1)=0,\qquad u(R_2)=0.
\]

令 \(k=\sqrt{2\mu E}/\hbar>0\)。一般解为
\[
u(r)=A\sin\big(k(r-R_1)\big)+B\cos\big(k(r-R_1)\big).
\]
由 \(u(R_1)=0\) 得 \(B=0\)，于是可以取
\[
u(r)=A\sin\big(k(r-R_1)\big).
\]
由 \(u(R_2)=0\) 得
\[
\sin\big(k(R_2-R_1)\big)=0 \quad\Longrightarrow\quad k_n=\frac{n\pi}{R_2-R_1},\quad n=1,2,\dots
\]
因此能量本征值为
\[
\boxed{\,E_n=\frac{\hbar^2 k_n^2}{2\mu}
=\frac{\hbar^2\pi^2 n^2}{2\mu\,(R_2-R_1)^2},\qquad n=1,2,\dots\,.}
\]

对应的径向函数（未归一）为
\[
u_n(r)=A_n\sin\!\Big(\frac{n\pi}{R_2-R_1}(r-R_1)\Big),\qquad R_1<r<R_2,
\]
并在区间外 \(u_n(r)=0\)。

归一化条件（球谐已归一化使得 \(\int|Y_{00}|^2\,\mathrm d\Omega=1\)）可写为
\[
\int_{R_1}^{R_2}|u_n(r)|^2\,\mathrm d r =1.
\]
令 \(L=R_2-R_1\)，做换元 \(x=r-R_1\) 得
\[
|A_n|^2\int_0^L\sin^2\!\Big(\frac{n\pi}{L}x\Big)\mathrm d x
=|A_n|^2\frac{L}{2}=1,
\]
于是
\[
A_n=\sqrt{\frac{2}{L}}=\sqrt{\frac{2}{R_2-R_1}}.
\]

因此归一化的径向波函数为
\[
\boxed{\,R_n(r)=\frac{u_n(r)}{r}
=\sqrt{\frac{2}{R_2-R_1}}\;\frac{1}{r}\,
\sin\!\Big(\frac{n\pi}{R_2-R_1}(r-R_1)\Big),\quad R_1<r<R_2\,.\,}
\]
完整的三维本征波函数为
\[
\boxed{\,\Psi_{n,\,l=0,\,m=0}(\boldsymbol r)
=\frac{1}{\sqrt{4\pi}}\,R_n(r)
=\frac{1}{\sqrt{4\pi}}\sqrt{\frac{2}{R_2-R_1}}\;
\frac{\sin\!\big(\dfrac{n\pi}{R_2-R_1}(r-R_1)\big)}{r},\quad R_1<r<R_2. \,}
\]
在 \(r\le R_1\) 或 \(r\ge R_2\) 时 \(\Psi=0\)。

以上即为 \(l=0\) 的能级与归一化本征波函数。
\end{solution}

\begin{problem}
  三维各向同性谐振子处于本征态 $\lvert Nlm\rangle$，求该状态下的 
$\langle r^2\rangle$、$\langle p^2\rangle$ 和 $\langle r^{-2}\rangle$。
提示：利用费曼–海尔曼定理和位力定理。
\end{problem}

\begin{solution}
  设三维各向同性谐振子哈密顿量为
\[
\hat H=\frac{\hat{\boldsymbol p}^2}{2\mu}+\frac{1}{2}\mu\omega^2\hat r^2,
\]
能量本征态用 $\lvert Nlm\rangle$ 表示（标准约定 $N=n_x+n_y+n_z=2n_r+l$），对应本征值
\[
E_{N}=\hbar\omega\Big(N+\frac{3}{2}\Big).
\]

\textbf{(1) \(\langle r^2\rangle\)}\quad 用费曼—海尔曼定理对参数 \(\omega\) 求导：
\[
\frac{\partial E_N}{\partial\omega}
=\Big\langle Nlm\Big|\frac{\partial\hat H}{\partial\omega}\Big|Nlm\Big\rangle
=\Big\langle Nlm\big|\mu\omega\hat r^2\big|Nlm\Big\rangle
=\mu\omega\langle r^2\rangle.
\]
由于 \(\partial E_N/\partial\omega=\hbar\big(N+\tfrac32\big)\)，得
\[
\boxed{\;\langle r^2\rangle=\frac{\hbar\,(N+\tfrac32)}{\mu\omega}\;}.
\]

\bigskip

\textbf{(2) \(\langle p^2\rangle\)}\quad 由谐振子的位力定理知道对二次势能有
\(\langle T\rangle=\langle V\rangle=E_N/2\)，且动能 \(\langle T\rangle=\langle \hat p^2\rangle/(2\mu)\)。于是
\[
\langle p^2\rangle=2\mu\langle T\rangle=\mu E_N
=\mu\hbar\omega\Big(N+\frac32\Big).
\]
因此
\[
\boxed{\;\langle p^2\rangle=\mu\hbar\omega\,(N+\tfrac32)\;},
\]


\bigskip

\textbf{(3) \(\langle r^{-2}\rangle\)}\quad 令参数化哈密顿量
\[
\hat H(\gamma)=\frac{\hat{\boldsymbol p}^2}{2\mu}+\frac{1}{2}\mu\omega^2\hat r^2
+\frac{\hbar^2\gamma}{2\mu}\frac{1}{\hat r^2},
\]
其中 \(\gamma\) 为连续参数；物理解出时取 \(\gamma=0\)。对 \(\gamma\) 应用费曼—海尔曼定理得
\[
\frac{\partial E_N(\gamma)}{\partial\gamma}
=\Big\langle Nlm;\gamma\Big|\frac{\partial\hat H(\gamma)}{\partial\gamma}\Big|Nlm;\gamma\Big\rangle
=\Big\langle Nlm;\gamma\Big|\frac{\hbar^2}{2\mu}\frac{1}{\hat r^2}\Big|Nlm;\gamma\Big\rangle.
\]
因此（在 \(\gamma=0\) 处）
\[
\langle r^{-2}\rangle=\frac{2\mu}{\hbar^2}\left.\frac{\partial E_N(\gamma)}{\partial\gamma}\right|_{\gamma=0}.
\]

下面给出被参量化问题的能级对 \(\gamma\) 的解析依赖（解径向方程可得，常见教材与文献可查）：
\[
E_N(\gamma)=\hbar\omega\Big(2n_r+1+\sqrt{\Big(l+\tfrac12\Big)^2+\gamma}\;\Big),
\]
这在 \(\gamma=0\) 时退化为 \(E_N=\hbar\omega(2n_r+l+\tfrac32)=\hbar\omega(N+\tfrac32)\)。于是
\[
\left.\frac{\partial E_N(\gamma)}{\partial\gamma}\right|_{\gamma=0}
=\hbar\omega\cdot\frac{1}{2\big(l+\tfrac12\big)}.
\]
代入上式得
\[
\langle r^{-2}\rangle=\frac{2\mu}{\hbar^2}\cdot\frac{\hbar\omega}{2\big(l+\tfrac12\big)}
=\frac{\mu\omega}{\hbar\big(l+\tfrac12\big)}.
\]
因此
\[
\boxed{\;\langle r^{-2}\rangle=\frac{\mu\omega}{\hbar\,(l+\tfrac12)}\; }.
\]
\end{solution}

\begin{problem}
  考虑三维各向同性的谐振子，

(1) 分别写出 $N=1$ 时，$(H_x, H_y, H_z)$ 表象和 $(H, l^2, l_z)$ 表象的三个基矢在球坐标系下的表达式。

(2) 验证变换矩阵的矩阵元  
$\langle n_r = 0,\, l = 1,\, m = 1 \vert n_x = 1,\, n_y = 1,\, n_z = 0 \rangle = -i/\sqrt{2}$。

\centerline{}
  
若粒子处于状态  
\[
\lvert \alpha \rangle = \dfrac{1}{2} \lvert 100 \rangle - \dfrac{i}{3} \lvert 010 \rangle + c \lvert 001 \rangle ,
\]
其中 $\lvert n_x, n_y, n_z \rangle$ 是表象 $\{H_x, H_y, H_z\}$ 下的基矢。

(3) 分别求 $E,\, l^2,\, l_z$ 的期望值；

(4) 写出其时刻 $t$ 的量子态。
\end{problem}

\begin{solution}
  设三维各向同性谐振子的单粒子基态能级记号为 $\lvert n_x,n_y,n_z\rangle$（各方向一维谐振子占据数），总体激发量子数
$N=n_x+n_y+n_z$。第一激发子空间 $N=1$ 的三个直角坐标基矢为
\[
\lvert x\rangle\equiv\lvert 100\rangle,\qquad
\lvert y\rangle\equiv\lvert 010\rangle,\qquad
\lvert z\rangle\equiv\lvert 001\rangle,
\]
其空间波函数分别与 $x,y,z$ 成正比：
\[
\psi_{100}(\boldsymbol r)\propto x\,e^{-\alpha^2 r^2/2},\quad
\psi_{010}(\boldsymbol r)\propto y\,e^{-\alpha^2 r^2/2},\quad
\psi_{001}(\boldsymbol r)\propto z\,e^{-\alpha^2 r^2/2},
\]
其中 $\alpha=\sqrt{\mu\omega/\hbar}$，$r,\theta,\phi$ 为球坐标。

(1) 在球坐标（即 $(H,l^2,l_z)$ 表象，对应 $n_r,l,m$）下，$N=1$ 对应唯一的径向量子数 $n_r=0$ 且 $l=1$（三重简并，$m=1,0,-1$）。利用
\[
x=r\sin\theta\cos\phi,\quad y=r\sin\theta\sin\phi,\quad z=r\cos\theta,
\]
以及球谐函数的线性组合，可写出两组基矢的角部分关系（忽略相同的径向因子 $r e^{-\alpha^2 r^2/2}$）：
\[
\begin{aligned}
&\text{直角坐标基（空间波函数形式）:}\\
&\quad \psi_x\propto r\sin\theta\cos\phi,\qquad
\psi_y\propto r\sin\theta\sin\phi,\qquad
\psi_z\propto r\cos\theta.
\end{aligned}
\]
而球谐基角部分为 $Y_{1m}(\theta,\phi)$，满足
\[
rY_{1,0}\propto r\cos\theta\propto z,
\]
并且
\[
rY_{1,\pm1}\propto \mp (x\pm i y).
\]
因此在态矢表示上可选取以下规范性的对应关系：
\[
\boxed{\;
\begin{aligned}
\lvert n_r=0,l=1,m=0\rangle&=\lvert z\rangle=\lvert 001\rangle,\\[4pt]
\lvert n_r=0,l=1,m=1\rangle&=-\frac{1}{\sqrt2}\big(\lvert x\rangle + i\lvert y\rangle\big)
=-\frac{1}{\sqrt2}\big(\lvert100\rangle + i\lvert010\rangle\big),\\[4pt]
\lvert n_r=0,l=1,m=-1\rangle&=\frac{1}{\sqrt2}\big(\lvert x\rangle - i\lvert y\rangle\big)
=\frac{1}{\sqrt2}\big(\lvert100\rangle - i\lvert010\rangle\big).
\end{aligned}
\; }
\]


(2) 题中需验证矩阵元
\[
\langle n_r=0,l=1,m=1\vert n_x=1,n_y=1,n_z=0\rangle=-\frac{\mathrm i}{\sqrt2}.
\]
验证是
\[
\langle n_r=0,l=1,m=1\vert n_x=0,n_y=1,n_z=0\rangle
=\langle 1,1\vert y\rangle.
\]
根据(1) 中的关系
\[
\lvert 1,1\rangle=-\frac{1}{\sqrt2}(\lvert x\rangle + i\lvert y\rangle),
\]
可得
\[
\langle 1,1\vert y\rangle=-\frac{i}{\sqrt2},
\]
即
\[
\boxed{\;
\langle n_r=0,l=1,m=1\vert n_x=0,n_y=1,n_z=0\rangle=-\frac{\mathrm i}{\sqrt2}.
\; }
\]

(3) 令
\[
\lvert\alpha\rangle=\frac{1}{2}\lvert100\rangle-\frac{\mathrm i}{3}\lvert010\rangle + c\lvert001\rangle,
\]
其中 $c$ 为复数常数。先计算能量的期望值。三维各向同性谐振子的能级与 $N=n_x+n_y+n_z$ 的关系为
\[
E(N)=\hbar\omega\Big(N+\frac{3}{2}\Big).
\]
对 $\lvert100\rangle,\lvert010\rangle,\lvert001\rangle$，均有 $N=1$，故它们的能量相同为 $E_1=\hbar\omega\big(1+\tfrac32\big)=\tfrac52\hbar\omega$。因此能量期望值为
\[
\boxed{\,\langle E\rangle=\frac{5}{2}\hbar\omega\,.}
\]

接着计算角动量平方与 $l_z$ 的期望值。由(1) 可见，三个直角坐标基矢都只含 $l=1$ 成分（$N=1$ 时唯一可能的 $l$ 值为 $1$），因此 $\lvert\alpha\rangle$ 是 $l=1$ 子空间的线性组合，故
\[
\langle l^2\rangle=\hbar^2 l(l+1)=\hbar^2\cdot 1\cdot2=2\hbar^2.
\]
对于 $\langle l_z\rangle$，把 $\lvert\alpha\rangle$ 在 $m$ 基上展开。利用(1) 的变换关系，得到 $m$ 分量振幅：
\[
\begin{aligned}
a_{+1}&=\langle 1,1\vert\alpha\rangle
=-\frac{1}{\sqrt2}\Big(\frac{1}{2}+ \mathrm i\cdot\big(-\frac{\mathrm i}{3}\big)\Big)
=-\frac{1}{\sqrt2}\Big(\frac{1}{2}+\frac{1}{3}\Big)
=-\frac{5}{6\sqrt2},\\[4pt]
a_{-1}&=\langle 1,-1\vert\alpha\rangle
=\frac{1}{\sqrt2}\Big(\frac{1}{2}+\mathrm i\cdot\big(-\frac{\mathrm i}{3}\big)\Big)
=\frac{5}{6\sqrt2},\\[4pt]
a_{0}&=\langle 1,0\vert\alpha\rangle=\langle z\vert\alpha\rangle=c.
\end{aligned}
\]
因此
\[
\langle l_z\rangle=\hbar\big( |a_{+1}|^2\cdot(+1)+|a_0|^2\cdot 0 + |a_{-1}|^2\cdot(-1)\big)
=\hbar\big(|a_{+1}|^2-|a_{-1}|^2\big)=0,
\]
因 $|a_{+1}|=|a_{-1}|$。综上：
\[
\boxed{\,\langle l^2\rangle=2\hbar^2,\qquad \langle l_z\rangle=0\,.}
\]

(4) 写出时刻 $t$ 的量子态。各基矢 $\lvert100\rangle,\lvert010\rangle,\lvert001\rangle$ 对应相同能量 $E_1=\tfrac52\hbar\omega$，因此在薛定谔绘景下整态仅得一个全局相位：
\[
\lvert\alpha(t)\rangle=e^{-\,\mathrm i E_1 t/\hbar}\lvert\alpha(0)\rangle
=e^{-\,\mathrm i\frac{5}{2}\omega t}\Big(\tfrac12\lvert100\rangle-\tfrac{\mathrm i}{3}\lvert010\rangle + c\lvert001\rangle\Big).
\]
若需写出坐标表示，只需在每一基矢前乘以相同因子 $e^{-i\frac52\omega t}$。
\end{solution}
\end{document}